\section{What the Guarantees Say}\label{sec:min-guarantees}

The theorem stack is conditional. It does not assert that arbitrary language
models preserve political meaning. It says what follows if the chosen state,
oracle, local laws, and sampling assumptions hold for the realized
application.

\paragraph{Root preservation.}
Fix a document partition and a binary merge tree with leaves in contiguous
order. Let $Z^{(1)}(x)$ be the stored root summary.

\begin{theorem}[Inductive preservation]\label{thm:one-pass}
If the realized run satisfies C1, C3, and
Assumption~\ref{ass:context}, then
\[
  \E\!\left[
    \metric\big(\fstar(Z^{(1)}(x)),\fstar(x)\big)
  \right]=0.
\]
\end{theorem}

The proof is the tree induction the picture suggests: C1 handles leaves,
context compatibility allows an oracle-equivalent child summary to be
substituted for the raw child span inside the parent, and C3 propagates the
claim through each merge.

\begin{theorem}[Preservation under re-summarization]\label{thm:multi-round}
If the run additionally satisfies C2, then the same equality holds for every
re-summary round $R\ge 1$ with $Z^{(R+1)}(x)=g(Z^{(R)}(x))$.
\end{theorem}

\begin{corollary}[Schedule invariance]\label{cor:schedule}
For a fixed contiguous partition, any two reduction schedules that satisfy
the same realized local laws have the same expected oracle value at the root.
\end{corollary}

This is the formal version of the leaf-size robustness claim. It does not say
every chunking is good; it says chunking and merge schedules are interchangeable
inside the class where the realized local-law checks hold.

\paragraph{Preference-objective equivalence.}
For DPO, GRPO, or related preference objectives, one more premise is needed:
the downstream preference must depend on the document only through the
preserved oracle value. Under that oracle-factorization assumption, training
on root summaries and training on full documents share the same population
target. In the approximate case, the deviation is controlled by the remaining
oracle distortion.

\begin{theorem}[Preference-objective equivalence;
\textsc{PaperPreferenceStack}]\label{thm:pref-equiv}
Under recompression plus oracle factorization and the application-specific
generator-indexing premises, zero residual gives identical full-document and
summary argmin sets. With residual $\varepsilon$, every minimizer of the
summary objective is $2\varepsilon$-optimal for the full objective.
\end{theorem}

\paragraph{Finite-sample certificate.}
Let
$\Delta_R=\E[\metric(\fstar(X),\fstar(Z^{(R)}(X)))]$ denote expected oracle
distortion after $R$ rounds.

\begin{theorem}[Unified preference gap]\label{thm:unified-gap}
If the method-specific expected inner loss has transport constant
$C_{\mathrm{meth}}$ in oracle distance, then
\[
  |G_{\mathrm{meth}}|\le C_{\mathrm{meth}}\Delta_R.
\]
\end{theorem}

\begin{theorem}[Audited preference-gap certificate;
\textsc{PaperErrorStack}]\label{thm:e2e}
Let $\hat\Delta_R$ be the design-based estimate of $\Delta_R$ from a
tree-unit audit with logged positive propensities. On the event where judge
calibration, sampling error, and clipping are bounded by
$B_{\mathrm{cal}}$, $B_{\mathrm{est}}$, and $B_{\mathrm{clip}}$ in objective
units,
\[
  |G_{\mathrm{meth}}|
  \le
  C_{\mathrm{meth}}\hat{\Delta}_R
  + B_{\mathrm{cal}} + B_{\mathrm{est}} + B_{\mathrm{clip}}.
\]
\end{theorem}

The local-law weight $\lambda$ in $J_\lambda$
(Section~\ref{sec:min-projection-estimation}, Eq.~\ref{eq:min-J-lambda})
feeds the certificate as a design choice: large $\lambda$ pulls
$\hat{\Delta}_R$ toward zero by projecting the operator into the
local-law subspace; small $\lambda$ leaves more residual for the audit
to report.

